%% abtex2-modelo-slides.tex, v-1.0 gfabinhomat
%% Copyright 2012-2015 by abnTeX2 group at http://www.abntex.net.br/ 
%%
%% This work may be distributed and/or modified under the
%% conditions of the LaTeX Project Public License, either version 1.3
%% of this license or (at your option) any later version.
%% The latest version of this license is in
%%   http://www.latex-project.org/lppl.txt
%% and version 1.3 or later is part of all distributions of LaTeX
%% version 2005/12/01 or later.
%%
%% This work has the LPPL maintenance status `maintained'.
%% 
%% The Current Maintainer of this work is Fábio Rodrigues Silva, 
%% member of abnTeX2 team, led by Lauro César Araujo. 
%% Further information are available on 
%% http://www.abntex.net.br/
%%
%% This work consists of the files abntex2-modelo-slides.tex, 
%% abntex2-modelo-references.bib and abntex2-modelo-marca.pdf
%%
%% Modelo desenvolvido por Fábio Rodrigues Silva (gfabinhomat@gmail.com)
%% Mais informações podem ser obtidas no guia do usuário Beamer 
%% (http://linorg.usp.br/CTAN/macros/latex/contrib/beamer/doc/beameruserguide.pdf)
%% Informações rápidas podem ser acessadas em http://en.wikibooks.org/wiki/LaTeX/Presentations


% Apresentações em widescreen. Outros valores possíveis: 1610, 149, 54, 43 e 32.
% Por padrão, as apresentações são no formato 4:3 (sem o aspectratio).
\documentclass[aspectratio=169]{beamer}	 	

\usetheme{Berlin}
\usecolortheme{dolphin}
\usefonttheme[onlymath]{serif}			% para fontes matemáticas
% Enconte mais temas e cores em http://www.hartwork.org/beamer-theme-matrix/ 
% Veja também http://deic.uab.es/~iblanes/beamer_gallery/index.html

% Customizações de Cores: fg significa cor do texto e bg é cor do fundo
% ---
% PACOTES
% ---
\usepackage[alf]{abntex2cite}		% Citações padrão ABNT
\usepackage[brazil]{babel}		% Idioma do documento
\usepackage{color}			% Controle das cores
\usepackage[T1]{fontenc}		% Selecao de codigos de fonte.
\usepackage{graphicx}			% Inclusão de gráficos
\usepackage[utf8]{inputenc}		% Codificacao do documento (conversão automática dos acentos)
\usepackage{txfonts}			% Fontes virtuais
% ---
\usepackage{caption}
% --- Informações do documento ---
\title{Adaptação de Segundo Nível, Parte 1}
\author{Pedro Yochinori Gushiken}
\institute{Universidade Federal do Rio Grande do Norte}
\date{\today}
% ---

% ----------------- INÍCIO DO DOCUMENTO --------------------------------------
\begin{document}

% ----------------- NOVO SLIDE --------------------------------
\begin{frame}

\begin{minipage}{1\linewidth}
  \centering
\end{minipage}

\titlepage

\end{frame}

% ----------------- NOVO SLIDE --------------------------------
\begin{frame}{Sumário}
\tableofcontents
\end{frame}

% ----------------- NOVO SLIDE --------------------------------
\section{Introdução}
\begin{frame}
\frametitle{Características}
\begin{enumerate}
\item<1-> Identificação, Caso Indireto\\ 
\item<2-> Propriedade do Fecho Convexo referente à Adaptação de Primeiro Nível\\
\item<3-> Projeção no Espaço de Parâmetros\\
\end{enumerate}
\end{frame}


% ----------------- NOVO SLIDE --------------------------------
\section{Conceitos}
\subsection{Problema do Controle Indireto Usando 1 Modelo de Identificação}

\begin{frame}
\frametitle{Equações de Primeiro Nível}




\centering
\begin{minipage}{.48\textwidth}
  \centering
\begin{equation}
\dot{x}_p(t) = A_p x_p(t) + b u(t)
\end{equation}

\begin{equation}
\dot{x}_m(t) = A_m x_m(t) + b r(t)
\end{equation}
\end{minipage}%
  \hfill
\begin{minipage}{.48\textwidth}
  \centering
\begin{equation}
\dot{\hat{x}}_p = A_m\hat{x}_p(t) + [A_1(t)-A_m]x_p(t) +bu(t)
\end{equation}
\begin{equation}
\theta_1 \rightarrow A_1 
\end{equation}

\end{minipage}



\end{frame}

\begin{frame}
\frametitle{Equações de Primeiro Nível, Cont.}

\centering
\begin{minipage}{.48\textwidth}
  \centering
\begin{equation}
u(t) = r(t) +k^T(t)x_p(t)
\end{equation}
\begin{equation}
k(t) = \theta_m - \theta_1 = k^* - \phi_1(t)
\end{equation}

\end{minipage}%
  \hfill
\begin{minipage}{.48\textwidth}
  \centering
\begin{equation}
e_1(t) = \hat{x}_p(t) - x_p(t)
\end{equation}
\begin{equation}
\dot{\theta_1} = \dot{\phi_1} = -{e_1}^TPbx_p(t)
\end{equation}

\end{minipage}

\end{frame}


\subsection{Múltiplos Modelos}
\begin{frame}
\frametitle{Múltiplos Modelos}

\begin{equation}
\Sigma_i: \dot{\hat{x}}_i = A_m\hat{x}_p(t) + [A_i(t)-A_m]x_p(t) +bu(t)
\end{equation}

\begin{equation}
\tilde{\theta} \equiv \phi
\end{equation}

\end{frame}

\begin{frame}

\centering
\begin{minipage}{.48\textwidth}
  \centering
\begin{equation}
\sum_{i=1}^{n+1}\alpha_i\theta_i(t_0) = \theta_p
\end{equation}
\begin{equation}
\dot{e}_i = A_m e_i(t) + b{\phi_i}^T(t)x_p(t)
\end{equation}
\begin{equation}
e_i(t_0)=0
\end{equation}
\end{minipage}%
  \hfill
\begin{minipage}{.48\textwidth}
  \centering

\begin{equation}
\dot{\theta}_i(t) = \dot{\phi}_i(t) = -{e_i}^T(t)Pbx_p(t)
\end{equation}
\begin{equation}
\theta_i(t_0) = \theta_{i0}
\end{equation}
\begin{equation}
\sum_{i=1}^{n+1}\alpha_i\phi(t_0)=0 
\end{equation}
\end{minipage}


\end{frame}


\begin{frame}

\centering
\begin{minipage}{.48\textwidth}
  \centering
  \begin{equation}
  \sum_{i=1}^{n+1}\alpha_i\theta_i(t) = \theta_p \Longleftrightarrow \bar{\Theta}(t)\bar{\alpha} = \theta_p \space 
  \end{equation}
\end{minipage}%
  \hfill

\begin{minipage}{.48\textwidth}
  \centering
    \begin{equation}
    \sum_{i=1}^{n+1}\alpha_i e_i(t) = 0 \Longleftrightarrow \bar{E}(t)\bar{\alpha} = 0 \space
    \end{equation}
\end{minipage}

\end{frame}

\begin{frame}

\begin{equation}
E(t)\alpha = -e_{n+1}(t)
\end{equation}

\begin{equation}
\dot{\hat{\alpha}} = E^T(t)E(t)\hat{\alpha} - E^T(t)e_{n+1}(t)  
\end{equation}

\end{frame}


% ----------------- NOVO SLIDE --------------------------------

% ----------------- NOVO SLIDE --------------------------------
\section{Referências}

% --- O comando \allowframebreaks ---
% Se o conteúdo não se encaixa em um quadro, a opção allowframebreaks instrui 
% beamer para quebrá-lo automaticamente entre dois ou mais quadros,
% mantendo o frametitle do primeiro quadro (dado como argumento) e acrescentando 
% um número romano ou algo parecido na continuação.
% --- O comando \allowframebreaks ---
% Se o conteúdo não se encaixa em um quadro, a opção allowframebreaks instrui 
% beamer para quebrá-lo automaticamente entre dois ou mais quadros,
% mantendo o frametitle do primeiro quadro (dado como argumento) e acrescentando 
% um número romano ou algo parecido na continuação.

\begin{frame}[allowframebreaks]
\frametitle{Referências}
\nocite{narendra1}
\nocite{narendra_main}
\nocite{narendra2012rationale}
\nocite{narendra1997}
\nocite{han1}
\nocite{han2010}
\nocite{han2012}
\nocite{kuipers1}
\nocite{ioannou2}
\nocite{chen}
\nocite{ogata}
\nocite{khalil}
\end{frame}
% ----------------- FIM DO DOCUMENTO -----------------------------------------
\bibliography{bib-controleadaptativo}
\end{document}
